\documentclass[../main.tex]{subfiles}

\begin{document}
\begin{CJK*}{UTF8}{gbsn}

\section*{Exercise 9}
Given data set $(X_1,Y_1),\ldots,(X_n,Y_n) \in \mathcal{X} \times \mathcal{Y}$
where $\mathcal{Y} = \mathbb{R}$,
consider split conformal prediction with 
two given models $\hat{f}_i: \mathcal{X} \to \mathbb{R}$, 
the score functions

\begin{equation*}
  s_i(x,y) = \abs{y - \hat{f}_i(x)},
\end{equation*}
and two split conformal prediction sets 

\begin{equation*}
\mathcal{C}_i(X_{n+1}) = \{y \in \mathcal{Y} \mid s_i(X_{n+1},y) \leqslant \hat{q}_i\},
\end{equation*}
where 

\begin{equation*}
    \hat{q}_i = \text{quantile}\curly{ \curly{s_i(X_j,Y_j)}_{j=1}^n ; (1-\alpha)(1+\frac{1}{n})}.
\end{equation*}
Put 

\begin{equation*}
    L_i = \frac{1}{n} \sum_{j=1}^n (Y_j-\hat{f}_i(X_j))^2
\end{equation*}
and $I = \argmin_{i \in \{1,2\}} L_i$.
Explain why 

\begin{equation*}
    P \curly{Y_{n+1} \in \mathcal{C}_I(X_{n+1})} \geqslant 1-\alpha
\end{equation*}
may not hold even if we assume that 
$(X_1,Y_1),\ldots,(X_{n+1},Y_{n+1})$
are exchangeable.

\paragraph{Solution}
In one sentence, the reason is that
even if we assume that the three random vectors 
$(X_1,Y_1),\ldots,(X_{n+1},Y_{n+1})$,
$s_1(X_1,Y_1) ,\ldots,s_1(X_{n+1},Y_{n+1})$,
and $s_2(X_1,Y_1) ,\ldots,s_2(X_{n+1},Y_{n+1})$
are all exchangeable, in general 
it is still not true that 
$s_I(X_1,Y_1) ,\ldots,s_I(X_{n+1},Y_{n+1})$
is exchangeable and thus 
the theorem guaranteeing the marginal coverage does not in general apply.

To show that things can indeed fail,
we give a simple counterexample.
Put $\mathcal{X} = \mathbb{R}^2$, $n = 1$, $\alpha = \frac{1}{2}$,
$X_j = (X_j^{(1)},X_j^{(2)})$, $X_j$ follows $N_2(0,I_2)$
distribution, $Y_j = X_j^{(1)} + X_j^{(2)}$,
$\hat{f}_1(x) = x^{(1)}$ and $\hat{f}_2(x) = x^{(2)}$.
It is clear that, for $i = 1,2$, $\hat{q}_i = s_i(X_1,Y_1)$,
$s_i(X_j, Y_j) = \abs{X_j^{(2-i)}}$,
and $L_i = (X_1^{(2-i)})^2$.
Thus,

\begin{equation*}
    \begin{split}
    & \tab[0.4cm] P \curly{s_I(X_2,Y_2) \leqslant \min(s_1(X_1,Y_1),s_2(X_1,Y_1))}
    \\ & = P \curly{s_1(X_2,Y_2) \leqslant \min(s_1(X_1,Y_1),s_2(X_1,Y_1)), L_1 < L_2}
    \\ & + P \curly{s_2(X_2,Y_2) \leqslant \min(s_1(X_1,Y_1),s_2(X_1,Y_1)), L_1 > L_2}
    \\ & = P \curly{\abs{X_2^{(2)}} \leqslant \abs{X_1^{(2)}} \leqslant \abs{X_1^{(1)}}}
    + P \curly{\abs{X_2^{(1)}} \leqslant \abs{X_1^{(1)}} \leqslant \abs{X_1^{(2)}}}
    = \frac{1}{6} + \frac{1}{6} = \frac{1}{3}
    \end{split}
\end{equation*}
and

\begin{equation*}
    P \curly{Y_2 \in \mathcal{C}_I(X_2)} 
    = P \curly{s_I(X_2,Y_2) \leqslant \min(s_1(X_1,Y_1),s_2(X_1,Y_1))}
    = \frac{1}{3} < \frac{1}{2} = 1-\alpha.
\end{equation*}
This completes the counterexample.
////
\end{CJK*}
\end{document}